\documentclass[aspectratio=1610, xcolor=table]{beamer}
% PACKAGES & CONFIGURATION
\usepackage[english]{babel}
\usepackage[utf8]{inputenc}
\usepackage{graphicx}
\usepackage{booktabs}
\usepackage{hyperref}
\usepackage{xcolor}
\usepackage{listings}
\usepackage{lstautogobble}
\usepackage{fancyvrb}

% THEME & COLORS
\mode<presentation>
{
\usetheme{default}
\usecolortheme{default}

% Define custom tech-forward color palette
\definecolor{PrimaryDark}{HTML}{1e293b}      % Slate 900 - Deep dark for backgrounds
\definecolor{PrimaryAccent}{HTML}{0ea5e9}    % Sky Blue - Modern accent
\definecolor{SecondaryAccent}{HTML}{10b981}  % Emerald - Secondary accent
\definecolor{TextLight}{HTML}{f1f5f9}        % Slate 100 - Light text
\definecolor{TextDark}{HTML}{1e293b}         % Slate 900 - Dark text
\definecolor{CodeBg}{HTML}{0f172a}           % Slate 950 - Code background
\definecolor{CodeText}{HTML}{e2e8f0}         % Slate 200 - Code text
\definecolor{AlertColor}{HTML}{f97316}       % Orange - For emphasis

% Set overall colors
\setbeamercolor{background canvas}{bg=PrimaryDark}
\setbeamercolor{normal text}{fg=TextLight}
\setbeamercolor{title}{fg=PrimaryAccent, bg=}
\setbeamercolor{subtitle}{fg=SecondaryAccent}
\setbeamercolor{structure}{fg=PrimaryAccent}
\setbeamercolor{alerted text}{fg=AlertColor}
\setbeamercolor{itemize item}{fg=PrimaryAccent}
\setbeamercolor{itemize subitem}{fg=SecondaryAccent}
\setbeamercolor{enumerate item}{fg=PrimaryAccent}
\setbeamercolor{enumerate subitem}{fg=SecondaryAccent}
\setbeamercolor{block title}{fg=TextLight, bg=PrimaryAccent}
\setbeamercolor{block body}{fg=TextLight, bg=}
\setbeamercolor{footline}{fg=TextLight, bg=}
\setbeamercolor{frametitle}{fg=PrimaryAccent}

% Set fonts
\usefonttheme{structurebold}
\setbeamerfont{frametitle}{size=\large, series=\bfseries}
\setbeamerfont{title}{size=\Huge, series=\bfseries}
\setbeamerfont{subtitle}{size=\Large}
\setbeamerfont{normal text}{size=\normalsize}
\setbeamerfont{itemize/enumerate body}{size=\large}
\setbeamerfont{itemize/enumerate subbody}{size=\normalsize}

% Customize templates
\setbeamertemplate{navigation symbols}{}  % Remove navigation symbols
\setbeamertemplate{caption}[numbered]

% Custom footer with slide number and author
\setbeamertemplate{footline}{
	\leavevmode%
	\hbox{%
		\begin{beamercolorbox}[wd=.5\paperwidth,ht=2.25ex,dp=1ex,left]{author in head/foot}%
			\hspace{1ex}\insertshortauthor
		\end{beamercolorbox}%
		\begin{beamercolorbox}[wd=.5\paperwidth,ht=2.25ex,dp=1ex,right]{date in head/foot}%
			\insertframenumber{}\hspace{1ex}
		\end{beamercolorbox}}%
	\vskip0pt%
}

% Customize frame title style
\setbeamertemplate{frametitle}{
\vskip0.5em
{\usebeamerfont{frametitle}\usebeamercolor[fg]{frametitle}\insertframetitle}
\vskip-0.5em
{\color{gray}\rule{\linewidth}{0.4pt}}
}
}

% CODE HIGHLIGHTING

% Terraform & HashiCorp Configuration Language (HCL)
\lstdefinelanguage{Terraform}{
	sensitive=true,
	morekeywords={
		terraform,required_providers,required_version,provider,resource,data,module,
		variable,locals,output,backend,cloud,workspace,lifecycle,depends_on,
		count,for_each,dynamic,provisioner,connection,validation,precondition,
		postcondition,check,assert,source,version,type,description,default,nullable,
		tags,bucket,region,rule,status,id,merge
	},
	morekeywords=[2]{true,false,null},
	morecomment=[l]{\#},
	morecomment=[l]{//},
	morecomment=[s]{/*}{*/},
	morestring=[b]"
}

\lstset{
	autogobble=true,
	basicstyle=\ttfamily\small\color{CodeText},
	backgroundcolor=\color{CodeBg},
	breaklines=true,
	captionpos=b,
	commentstyle=\color{gray},
	keywordstyle=\color{PrimaryAccent}\bfseries,
	numberstyle=\tiny\color{gray},
	stringstyle=\color{SecondaryAccent},
	showstringspaces=false,
	frame=single,
	frameround=tttt,
	framerule=0pt,
	rulecolor=\color{gray},
	xleftmargin=1em,
	xrightmargin=1em,
	framexleftmargin=0.8em,
	framexrightmargin=0.8em,
}


\title{Measure First, Optimise Later}
\subtitle{Benchmarking in C\#}
\author{Marcel Barlik}
\date{\today}


\begin{document}

\begin{frame}[plain]
	\titlepage
\end{frame}

\begin{frame}{Agenda}
	\begin{itemize}
		\item \textbf{Context \& Motivation} --- Why performance measurement matters
		\item \textbf{What is Benchmarking?} --- Definitions and what we can measure
		\item \textbf{Key Benefits} --- What you gain by benchmarking
		\item \textbf{BenchmarkDotNet} --- The standard tool for C\#
		\item \textbf{Key Takeaways} --- What to do next
	\end{itemize}
	\vspace{1.5em}
	{\small \textcolor{gray}{Expected time: 10-15 minutes}}
\end{frame}


\section{Context \& Motivation}

\begin{frame}{The Performance Problem}
	\begin{columns}[T]
		\column{0.5\textwidth}
		\textbf{Symptoms we hear}
		\begin{itemize}
			\item ``This function feels slow''
			\item ``P99 latencies are out of control''
			\item ``We are spending a lot on cloud compute''
			\item Users complain about long app loading times
		\end{itemize}

		\column{0.5\textwidth}
		\textbf{Why they're hard to fix}
		\begin{itemize}
			\item ``Slow'' is not actionable
			\item We don't know \textit{where} the time goes
			\item Fixes are guesswork without data
			\item Regressions go unnoticed until production
		\end{itemize}
	\end{columns}

	\vspace{1.5em}
	\begin{alertblock}{The core problem}
		You cannot optimize what you have not measured.
	\end{alertblock}
\end{frame}

\begin{frame}{The Chain Is Only as Strong as Its Weakest Link}
	A service can only be as fast as the slowest dependency it relies on.

	\vspace{1em}
	\textbf{Common bottleneck dependencies}
	\begin{itemize}
		\item Database queries and round-trips
		\item External API calls
		\item Message queues and event buses
		\item Cache misses (falling back to slower sources)
		\item Network I/O and serialisation
	\end{itemize}

	\vspace{0.5em}
	\textcolor{SecondaryAccent}{\textbf{The analogy:}}
	``A Ferrari on a pebble road drives at pebble-road speed.''

	\vspace{0.5em}
	\begin{alertblock}{Key Insight}
		Identify your actual bottlenecks before optimising.
		Often it's not the code -- it's what the code is waiting for.
	\end{alertblock}
\end{frame}

\begin{frame}{Intuition vs. Reality}
	\begin{columns}[T]
		\column{0.5\textwidth}
		\textbf{What developers assume is slow}
		\begin{itemize}
			\item String operations
			\item Simple loops
			\item Math calculations
			\item Reading from memory
		\end{itemize}

		\column{0.5\textwidth}
		\textbf{What profilers actually find}
		\begin{itemize}
			\item Database round-trips
			\item Network and I/O calls
			\item Serialisation overhead
			\item Lock contention and GC pressure
		\end{itemize}
	\end{columns}

	\vspace{1.5em}
	\begin{block}{Donald Knuth, 1974}
		``We should forget about small efficiencies, say about 97\% of the time: premature optimization is the root of all evil.
		Yet we should not pass up our opportunities in that critical 3\%.''
	\end{block}
\end{frame}


\section{What is Benchmarking?}

\begin{frame}{What is Benchmarking?}
	\begin{block}{Definition}
		Benchmarking is the practice of \textbf{measuring the performance of a
			specific piece of code} under controlled, repeatable conditions -- producing
		data you can compare, track, and act on.
	\end{block}

	\vspace{1em}
	\textbf{What makes a good benchmark?}
	\begin{itemize}
		\item \textbf{Isolated} -- tests one thing at a time
		\item \textbf{Repeatable} -- same code produces comparable results
		\item \textbf{Statistically sound} -- accounts for warmup, JIT, and noise
		\item \textbf{Meaningful} -- measures something that reflects real usage
	\end{itemize}
\end{frame}


\begin{frame}{Micro vs.\ Macro Benchmarks}
	\begin{columns}[T]
		\column{0.5\textwidth}
		\textbf{Micro Benchmarks}
		\begin{itemize}
			\item Measure a single method
			\item Fast to run, easy to iterate
			\item Tool: \texttt{BenchmarkDotNet}
		\end{itemize}
		\vspace{0.5em}
		{\small \textcolor{SecondaryAccent}{Best for:} comparing implementations,
			validating a specific optimisation, tracking regressions at the unit level}

		\column{0.5\textwidth}
		\textbf{Macro Benchmarks}
		\begin{itemize}
			\item Measure end-to-end behaviour
			\item Closer to real-world load
			\item Tools: load testing, profilers, APM
		\end{itemize}
		\vspace{0.5em}
		{\small \textcolor{SecondaryAccent}{Best for:} finding real bottlenecks in
			production scenarios, validating system-level performance goals}
	\end{columns}

	\vspace{1em}
	\textcolor{PrimaryAccent}{\textbf{Both are valuable}} -- they answer different questions.
	Start with micro when optimising a specific hot path.
\end{frame}

\begin{frame}{What Can We Measure?}
	\begin{columns}[T]
		\column{0.5\textwidth}
		\textbf{Time}
		\begin{itemize}
			\item Mean execution time (ns / \textmu{}s / ms)
			\item Min, max, percentiles
			\item Statistical error and variance
		\end{itemize}

		\vspace{0.8em}
		\textbf{Memory}
		\begin{itemize}
			\item Heap allocations per operation
			\item GC collection frequency
			\item Peak working set
		\end{itemize}

		\column{0.5\textwidth}
		\textbf{Throughput}
		\begin{itemize}
			\item Operations per second
			\item Requests per second
			\item Bytes processed per second
		\end{itemize}

		\vspace{0.8em}
		\textbf{Hardware counters}
		\begin{itemize}
			\item CPU cache misses
			\item Branch mispredictions
			\item CPU cycles and instructions
		\end{itemize}
	\end{columns}
\end{frame}


\section{Key Benefits}

\begin{frame}{Without vs.\ With Benchmarks}
	\begin{columns}[T]
		\column{0.5\textwidth}
		\textbf{Without Benchmarks}
		\begin{itemize}
			\item Optimise based on gut feeling
			\item Cannot prove a fix actually helped
			\item Regressions slip into production unnoticed
			\item Performance reviews become debates
		\end{itemize}

		\column{0.5\textwidth}
		\textbf{With Benchmarks}
		\begin{itemize}
			\item Optimise based on evidence
			\item Prove improvements with numbers
			\item Catch regressions in CI before they ship
			\item Performance reviews become conversations
		\end{itemize}
	\end{columns}
\end{frame}

\begin{frame}{Key Benefits of Benchmarking}
	\begin{block}{Data-driven Decisions}
		Optimise the code that is \textit{actually} slow, not the code you assume is slow.
	\end{block}

	\vspace{0.6em}

	\begin{block}{Catch Regressions Early}
		Benchmarks can run in CI. A performance regression often indicates a bug or oversight.
	\end{block}

	\vspace{0.6em}

	\begin{block}{Validate Improvements}
		Quantifiable wins and side-by-side comparissons.
	\end{block}

	\vspace{0.6em}

	\begin{block}{Compare Alternatives Confidently}
		Choosing between two implementations? The numbers help decide.
	\end{block}
\end{frame}

\begin{frame}{Where Benchmarking Fits in the Dev Lifecycle}
	\begin{enumerate}
		\item \textbf{Development} -- Benchmark a hot path while writing or
		      refactoring it. Fail fast, iterate quickly.

		      \vspace{0.5em}
		\item \textbf{Code Review} -- Attach benchmark results to the pull request
		      as evidence. Performance becomes part of the review.

		      \vspace{0.5em}
		\item \textbf{CI / CD} -- Run a benchmark suite automatically.
		      Gate merges on regressions beyond an agreed threshold.

		      \vspace{0.5em}
		\item \textbf{Production} -- Profile real workloads to discover the next
		      bottleneck and feed that back into step 1.
	\end{enumerate}

	\vspace{1em}
	{\textcolor{SecondaryAccent}{\textbf{Key insight:}}}
	Performance is a feature -- treat it like one throughout the whole cycle.
\end{frame}


\section{BenchmarkDotNet}

\begin{frame}{BenchmarkDotNet}
	\textbf{The standard micro-benchmarking library for .NET}

	\vspace{0.8em}
	\begin{columns}[T]
		\column{0.6\textwidth}
		\textbf{What it handles for you}
		\begin{itemize}
			\item JIT warmup and tiered compilation
			\item Multiple iterations and statistical analysis
			\item Outlier detection and error reporting
			\item Memory allocation tracking (\texttt{[MemoryDiagnoser]})
			\item Baseline comparisons and ratios
			\item Export to Markdown, HTML, CSV
		\end{itemize}

		\column{0.4\textwidth}
		\textbf{Credibility}
		\begin{itemize}
			\item Used by the .NET runtime team
			\item Used by Microsoft ASP.NET team
			\item Industry standard across the .NET ecosystem
		\end{itemize}
	\end{columns}

	\vspace{1em}
	\textbf{Get started in one command:}\\
	\vspace{0.3em}
	\texttt{\textcolor{SecondaryAccent}{dotnet add package BenchmarkDotNet}}
\end{frame}

\begin{frame}[fragile]{A Concrete C\# Example}
	\textbf{Comparing string concatenation approaches}

	\begin{lstlisting}[language={[Sharp]C}, basicstyle=\ttfamily\tiny\color{CodeText}, numbers=left]
[MemoryDiagnoser]
public class StringBenchmarks
{
    [Params(10, 100, 1000)]
    public int N;

    [Benchmark(Baseline = true)]
    public string Concatenation()
    {
        var s = "";
        for (var i = 0; i < N; i++) s += i;
        return s;
    }

    [Benchmark]
    public string UseStringBuilder()
    {
        var sb = new StringBuilder();
        for (var i = 0; i < N; i++) sb.Append(i);
        return sb.ToString();
    }
}
	\end{lstlisting}
\end{frame}

\begin{frame}{Reading the Results}
	\textbf{Real results -- N = 1000} {\small \textcolor{PrimaryAccent}{\hyperlink{full-results}{(full results in appendix})}}

	\vspace{0.8em}
	\begin{center}
		\small
		\begin{tabular}{l|r|r|r|r|r|r}
			\toprule
			\textbf{Method}  & \textbf{Mean}    & \textbf{Error}  & \textbf{StdDev} & \textbf{Ratio} & \textbf{Gen0} & \textbf{Allocated} \\
			\midrule
			Concatenation    & 56.45 \textmu{}s & 0.31 \textmu{}s & 0.24 \textmu{}s & 1.00           & 56.58         & 2,774 KB           \\
			UseStringBuilder & 2.16 \textmu{}s  & 0.01 \textmu{}s & 0.01 \textmu{}s & 0.04           & 0.29          & 14.4 KB            \\
			\bottomrule
		\end{tabular}
	\end{center}

	\vspace{1em}
	\begin{columns}[T]
		\column{0.5\textwidth}
		\begin{itemize}
			\item \textbf{Mean} -- average execution time per operation
			\item \textbf{Error / StdDev} -- statistical noise in the measurement
		\end{itemize}

		\column{0.5\textwidth}
		\begin{itemize}
			\item \textbf{Ratio} -- relative to the baseline (1.00)
			\item \textbf{Gen0} -- GC collections per 1,000 operations
			\item \textbf{Allocated} -- heap memory per operation
		\end{itemize}
	\end{columns}

	\vspace{0.8em}
	\textcolor{SecondaryAccent}{\textbf{Result:}}
	\texttt{StringBuilder} is \textbf{25$\times$ faster}, triggers \textbf{195$\times$ fewer GC collections}, and allocates \textbf{192$\times$ less memory}.
\end{frame}


\begin{frame}{BenchmarkDotNet: Do's and Don'ts}
	\begin{columns}[T]
		\column{0.5\textwidth}
		\textbf{\textcolor{green!70!black}{Do's}}
		\begin{itemize}
			\item Use \texttt{[MemoryDiagnoser]} to track allocations
			\item Use \texttt{[Params(...)]} to test multiple scenarios
			\item Compare apples-to-apples -- same inputs, same environment
			\item When running on UNIX-like systems, use \texttt{nice -n -19} to reduce noise from other processeses
		\end{itemize}

		\column{0.5\textwidth}
		\textbf{\textcolor{red!70!black}{Don'ts}}
		\begin{itemize}
			\item Don't benchmark in Debug mode
			\item Don't include I/O or network calls in micro-benchmarks
			\item Don't optimize without data
		\end{itemize}
	\end{columns}
\end{frame}


\section{Closing}

\begin{frame}{Key Takeaways}
	\begin{enumerate}
		\item \textbf{You cannot optimise what you have not measured} --
		      vague complaints are not actionable; numbers are

		\item \textbf{Developer intuition about bottlenecks is often wrong} --
		      let a profiler or benchmark find the real hot path

		\item \textbf{Benchmarking turns assumptions into evidence} --
		      before/after numbers end debates and build confidence

		\item \textbf{BenchmarkDotNet is the standard tool for .NET} --
		      it handles the hard parts (warmup, statistics, allocations)

		\item \textbf{Integrate early} --
		      benchmarks in CI catch regressions before they reach production
	\end{enumerate}

	\vspace{1em}
	{\large \textcolor{PrimaryAccent}{\textbf{Next step:}}}
	Look through your most hit service, and find one hot path to benchmark and optimise.
\end{frame}

\begin{frame}{Questions \& Discussion}
	\begin{center}
		\vspace{2em}
		\Huge \textcolor{PrimaryAccent}{Questions?}
		\vspace{2em}

		\normalsize
		What performance problems are you dealing with right now?

		\vspace{1em}
		{\small Have you already used BenchmarkDotNet? What did you learn?}
	\end{center}
\end{frame}


\appendix

\begin{frame}[allowframebreaks]{References \& Further Reading}
	\small
	\begin{itemize}
		\item \textbf{BenchmarkDotNet} -- \url{https://benchmarkdotnet.org}
		      -- official docs, guides, and configuration reference

		\item \textbf{BenchmarkDotNet GitHub} --
		      \url{https://github.com/dotnet/BenchmarkDotNet} --
		      source, issues, and real-world usage examples from the .NET ecosystem

		\item \textbf{.NET Blog: Performance Improvements} --
		      \url{https://devblogs.microsoft.com/dotnet/} --
		      the .NET team publishes benchmark-driven perf posts each release

        \item \textbf{.NET Blog: .NET 10 Improvements} --
              \url{https://devblogs.microsoft.com/dotnet/performance-improvements-in-net-10/} --
              detailed benchmarks showing the impact of .NET 10 optimisations
	\end{itemize}
\end{frame}

\begin{frame}[fragile]{\hypertarget{full-results}{Benchmark Environment}}
	\begin{lstlisting}[language={}, basicstyle=\ttfamily\scriptsize\color{CodeText}, autogobble=false]
BenchmarkDotNet v0.15.8, Linux NixOS 25.11 (Xantusia)
AMD Ryzen 7 9850X3D 5.61GHz, 1 CPU, 16 logical and 8 physical cores
.NET SDK 10.0.201
[Host]     : .NET 10.0.5 (10.0.5, 10.0.526.15411), X64 RyuJIT x86-64-v4
DefaultJob : .NET 10.0.5 (10.0.5, 10.0.526.15411), X64 RyuJIT x86-64-v4
    \end{lstlisting}

	\vspace{1.6em}

	\resizebox{\textwidth}{!}{%
		\begin{tabular}{l r r r r r r r r r r}
			\toprule
			Method           & N    & Mean        & Error      & StdDev     & Ratio & RatioSD & Gen0    & Gen1   & Allocated & Alloc Ratio \\
			\midrule
			Concatenation    & 10   & 53.37 ns    & 0.690 ns   & 0.646 ns   & 1.00  & 0.02    & 0.0067  & --     & 336 B     & 1.00        \\
			UseStringBuilder & 10   & 24.01 ns    & 0.180 ns   & 0.159 ns   & 0.45  & 0.01    & 0.0030  & --     & 152 B     & 0.45        \\
			\addlinespace
			Concatenation    & 100  & 948.49 ns   & 5.435 ns   & 4.818 ns   & 1.00  & 0.01    & 0.4139  & --     & 20856 B   & 1.00        \\
			UseStringBuilder & 100  & 228.48 ns   & 3.308 ns   & 3.094 ns   & 0.24  & 0.00    & 0.0253  & --     & 1280 B    & 0.06        \\
			\addlinespace
			Concatenation    & 1000 & 56452.99 ns & 310.789 ns & 242.644 ns & 1.00  & 0.01    & 56.5796 & 0.7935 & 2840456 B & 1.000       \\
			UseStringBuilder & 1000 & 2163.83 ns  & 12.941 ns  & 10.104 ns  & 0.04  & 0.00    & 0.2899  & 0.0076 & 14712 B   & 0.005       \\
			\bottomrule
		\end{tabular}}
\end{frame}

\end{document}
